\documentclass[11pt]{amsart}
\usepackage{geometry} % See geometry.pdf to learn the layout options. There are lots.
\geometry{letterpaper} % ... or a4paper or a5paper or ...
%\geometry{landscape} % Activate for for rotated page geometry
%\usepackage[parfill]{parskip} % Activate to begin paragraphs with an empty line rather than an indent
\usepackage{graphicx}
\usepackage{amssymb}
\usepackage{epstopdf}
\usepackage{multirow}

\DeclareGraphicsRule{.tif}{png}{.png}{convert #1 dirname #1/basename #1 .tif`.png}

\title{Ejemplo 2.6}
%\date{}
% Activate to display a given date or no date

\begin{document}
\maketitle
%\section{}
%\subsection{}
\noindent Analicemos el siguiente caso.

\noindent Nota: 
\begin{quotation}
\noindent Observe el lector que si ei selector toma el valor l o Z se realiza la acción 1, si el selector toma el valor 3, 4 o 5 se realiza la acción 2, y si el selector toma cualquier otro valor se realiza la acción 3. Luego, cuando se continúa con el flujo normal del diagrama se realiza la acción X.
\end{quotation}

\noindent A continuación presentamos al diagrama de flujo 2.9 en lenguaje algorítmico. A continuación presentamos algunos ejemplos donde el lector puede aplicar los conceptos estudiados con la estructura selectiva si múltiple.




\end{document}